% !TEX program = xelatex
% 数学物理方法1 · 第三章专题 · 级数乘法与 Cauchy 乘积
\documentclass[11pt]{ctexart}
%=====================================================================
% theme-violet.tex —— 讲义精讲 · 紫花主题
% 用法：\documentclass[11pt]{ctexart}
%       \input{assets/theme-violet}
%       \begin{document} ... \end{document}
% 只用 ctexart 自带能力 + xcolor/enumitem/fancyhdr/pifont/amsmath，
% 不依赖 tikz/tcolorbox，保证 MiKTeX 可直接编译。
%=====================================================================

\usepackage[a4paper,left=2.0cm,right=2.0cm,top=2.1cm,bottom=2.3cm,%
            headsep=12pt,footskip=22pt]{geometry}
\usepackage{amsmath,amssymb}
\usepackage{needspace}
\usepackage{xcolor}
\usepackage{array}
\usepackage{enumitem}
\usepackage{pifont}
\usepackage{fancyhdr}
\usepackage[hidelinks]{hyperref}

%---------------------------------------------------------------- 配色
\definecolor{violet}{RGB}{107,78,143}      % 主色：紫
\definecolor{violetdeep}{RGB}{63,44,92}    % 深紫：标题/强调
\definecolor{petal}{RGB}{183,159,208}      % 花瓣：装饰线、花饰
\definecolor{bandviolet}{RGB}{247,243,251} % 浅紫底：档案/提示框
\definecolor{bandcream}{RGB}{252,249,242}  % 米白底：易错/注意框
\definecolor{rulepurple}{RGB}{217,207,230} % 分隔线
\definecolor{inkgray}{RGB}{92,96,110}      % 次要文字

%---------------------------------------------------------------- 字体
% 有 Noto CJK 用 Noto，没有就退回 ctex 默认（Windows 下为宋体/黑体）
\IfFontExistsTF{Noto Serif CJK SC}{\setCJKmainfont{Noto Serif CJK SC}}{}
\IfFontExistsTF{Noto Sans CJK SC}{\setCJKsansfont{Noto Sans CJK SC}}{}

%---------------------------------------------------------------- 版式
\setlist{nosep,leftmargin=1.9em}
\linespread{1.28}
\setlength{\parskip}{3.5pt}
\setlength{\parindent}{0pt}
\setlength{\tabcolsep}{5pt}
\renewcommand{\arraystretch}{1.28}

%---------------------------------------------------------------- 花饰
% \ding{95} = 五瓣小花，\ding{96} = 星形花；都来自 pifont，必定存在
\newcommand{\orn}{\textcolor{petal}{\ding{95}}}          % 行内小标
\newcommand{\ornstar}{\textcolor{petal}{\ding{96}}}      % 星形花
\newcommand{\flowerrule}{%                                % 居中花饰分隔线
  \par\vspace{7pt}%
  \noindent\makebox[\linewidth]{\color{petal}%
    \rule{0.20\linewidth}{0.6pt}\hspace{8pt}{\small\ding{95}}\hspace{4pt}%
    \raisebox{1.1pt}{\tiny\ding{95}}\hspace{4pt}{\small\ding{95}}\hspace{8pt}%
    \rule{0.20\linewidth}{0.6pt}}%
  \par\vspace{7pt}}

%---------------------------------------------------------------- 标题与节头
\newcommand{\doctitle}[3]{%   #1 主标题  #2 副标题  #3 元信息（可用 \\ 分行）
  \begin{center}
    {\color{petal}\small \orn\enspace #2\enspace \ornstar}\\[7pt]
    {\LARGE\bfseries\color{violetdeep} #1}\\[10pt]
    {\color{petal}\rule{0.55\linewidth}{1pt}}\\[4pt]
    {\small\color{inkgray} #3}
  \end{center}\vspace{2pt}}

\newcommand{\hdr}[2]{%   #1 章节号（如 §0）  #2 标题
  \par\vspace{12pt}%
  \noindent\colorbox{violetdeep}{\color{white}\bfseries\small\;#1\;}\hspace{7pt}%
  {\color{violetdeep}\large\bfseries #2}\par\nobreak
  \vspace{3pt}{\color{petal}\hrule height 1.2pt}\vspace{6pt}}

\newcommand{\subhdr}[1]{%
  \par\vspace{6pt}%
  \noindent{\color{violet}\bfseries #1}\par\vspace{1pt}}

%---------------------------------------------------------------- 提示框
\newcommand{\callout}[1]{% 浅紫底：讲义档案、读法说明
  \par\medskip\noindent
  \fcolorbox{rulepurple}{bandviolet}{\parbox{\dimexpr\linewidth-2\fboxsep-2\fboxrule}{\vspace{2pt}#1\vspace{2pt}}}%
  \par\medskip}

\newcommand{\warnbox}[1]{% 米白底：易错、边界条件
  \par\medskip\noindent
  \fcolorbox{rulepurple}{bandcream}{\parbox{\dimexpr\linewidth-2\fboxsep-2\fboxrule}{\vspace{2pt}#1\vspace{2pt}}}%
  \par\medskip}

\newcommand{\gapitem}[2]{% #1 缺口类型（紫） #2 内容
  \par\vspace{2pt}\noindent
  \hangindent=1.4em\hangafter=1
  {\color{petal}\small\ding{96}}\enspace{\bfseries\color{violetdeep}#1}\enspace #2\par}

% 题干块（\qstem）：每道例题／作业的「详解」开头必须先复述题干与出处。
% 起因：2026-10-25 第五章讲义把七道作业的题面集中放在一页、详解另起一页，
% 读者翻到「第 4 题」时看不到条件，等于一道没有题干的题目。
% 规则：详解块内自带题干，读者不翻页就能读懂题在问什么。
\newcommand{\qstem}[1]{%
  \par\medskip\noindent
  \fcolorbox{violet}{bandviolet}{%
    \parbox{\dimexpr\linewidth-2\fboxsep-2\fboxrule}{%
      \vspace{2pt}%
      {\small\bfseries\color{violet}题干}\hspace{0.8em}#1%
      \vspace{2pt}}}%
  \par\medskip}

% v2: 保护题干与开头，长题可跨页，续页补上下文。
\newcommand{\qbreak}{\needspace{6\baselineskip}}
% 独立练习：ID、来源、题干、考点、答案标签（通常 ans:ID）。
\newcommand{\practice}[5]{%
  \qbreak\subhdr{练习 #1}\qstem{（#2）#3}%
  \textbf{考点}：#4\par
  详解见附录第~\pageref{#5}~页。\par\medskip}
% 附录开头：ID、来源、条件摘要；自动产生 ans:ID 标签。
\newcommand{\answerstart}[3]{%
  \qbreak\subhdr{#1\quad 详解}\label{ans:#1}\qstem{（#2）#3}}
\newcommand{\answercontinue}[2]{%
  \needspace{5\baselineskip}\subhdr{#1\quad 续解}%
  {\small\color{inkgray}本页接续：#2}\par\medskip}
% 延迟提供默认值，让旧讲义可在导言区保留自己的 \newcommand{\step}。
\AtBeginDocument{\providecommand{\step}[1]{\needspace{3\baselineskip}\par\medskip\textbf{#1}\par}}

% 前提旁注：只放在「需要理解的那一步」后面，给陈述 + 本题代入，不重证。
\newcommand{\aside}[1]{%
  \par\vspace{1pt}\noindent
  \hangindent=1.6em\hangafter=1
  {\small\color{inkgray}%
    \hspace*{0.6em}{\color{petal}$\triangleright$}\hspace{2pt}#1}%
  \par\vspace{2pt}}

%---------------------------------------------------------------- 表格
% 本机 MiKTeX 的 colortbl 与 array 版本不匹配（\rowcolor 会报错），
% 故表格统一用「上粗紫线 + 表头 + 中细线」的三线式，不依赖 colortbl。
\newcommand{\tabtop}{\par\vspace{3pt}{\color{violet}\hrule height 1.2pt}\vspace{1pt}}
\newcommand{\tabmid}{\par\vspace{1pt}{\color{petal}\hrule height 0.5pt}\vspace{1pt}}
\newcommand{\tabbot}{\par\vspace{3pt}{\color{rulepurple}\hrule height 0.5pt}\vspace{8pt}}
\newcommand{\thead}[1]{{\bfseries\color{violetdeep}#1}}

%---------------------------------------------------------------- 页眉页脚
\pagestyle{fancy}
\fancyhf{}
\renewcommand{\headrulewidth}{0.4pt}
\renewcommand{\headrule}{\hbox to\headwidth{\color{rulepurple}\leaders\hrule height \headrulewidth\hfill}}
\fancyhead[L]{\small\color{inkgray}\songti \doctag}
\fancyhead[R]{\small\color{inkgray}\songti \coursetag}
\fancyfoot[C]{\small\color{inkgray}\orn\enspace{\color{violet}\thepage}\enspace\orn}

% 每篇讲义在导言区设这两个值
\newcommand{\doctag}{}
\newcommand{\coursetag}{}
\newcommand{\settags}[2]{\renewcommand{\doctag}{#1}\renewcommand{\coursetag}{#2}}

%---------------------------------------------------------------- 数学记号
\newcommand{\vect}[1]{\boldsymbol{#1}}
\newcommand{\dd}{\,\mathrm{d}}
\newcommand{\ddx}{\frac{\mathrm{d}}{\mathrm{d}x}}
\addtolength{\jot}{2pt}



\settags{第三章专题 · 级数乘法与 Cauchy 乘积}{数学物理方法1}

\renewcommand{\qbreak}{\needspace{0.45\textheight}}
\let\origsubhdr\subhdr
\renewcommand{\subhdr}[1]{\par\needspace{4\baselineskip}\origsubhdr{#1}}
\let\orighdr\hdr
\renewcommand{\hdr}[2]{\par\needspace{9\baselineskip}\orighdr{#1}{#2}}
\setlength{\abovedisplayskip}{7pt plus 2pt minus 2pt}
\setlength{\belowdisplayskip}{7pt plus 2pt minus 2pt}

\begin{document}

\doctitle{第三章专题　级数乘法与 Cauchy 乘积}
{为什么是“指数相加”、什么时候才能乘、双边级数怎么配对}
{依据课件《第三章 解析函数的幂级数展开》p19--p20、p45--p47\quad 2026-11-01\quad 数学物理方法1}

\hdr{§0}{这个专题解决哪三个问题}

课件 p19 写了一行“幂级数可以相乘”，p45--p47 的贝塞尔例题又要求把两个指数级数乘起来，但没有把中间那一步讲透。这份专题只回答三个问题：

\begin{enumerate}
\item 为什么两个级数相乘时，系数要按\textbf{指数相加}来配，而不是按“下标相加”？
\item 什么时候可以放心地把两个无限和乘开、再重新分组？条件收敛行不行？
\item 遇到洛朗级数（有负幂项）时，配对的规则要不要改？
\end{enumerate}

这三个问题分别对应 §1、§2、§3；§4 给出三条可操作的算法，§5 用五道例题把它们跑一遍，§6 列易错点，§7 留六道自检题（答案在最后）。

\aside{与第三章讲义的关系：讲义 §2.4 只写了乘法公式的结论，例 9（贝塞尔）用了 $k-m=n$ 的配对。本专题把这两处的“为什么”补上，不重复讲幂级数的收敛半径。}

\hdr{§1}{Cauchy 乘积：指数相加，不是下标相加}

\subhdr{1.1\quad 先看两个多项式}

设

\[ A(z)=a_0+a_1z+a_2z^2,\qquad B(z)=b_0+b_1z+b_2z^2 . \]

把它们乘开，每一项都是“从第一行取一个单项、从第二行取一个单项”：

\[ A(z)B(z)=\sum_{k=0}^{2}\sum_{m=0}^{2}a_kb_mz^{k+m} . \]

要收集 $z^n$ 的系数，就得挑出所有 $k+m=n$ 的项：

\[ c_n=\sum_{k=0}^{n}a_kb_{n-k}\qquad(n=0,1,2,3,4) . \]

注意这里 $a_k$ 的“指数”是 $k$，$b_m$ 的“指数”是 $m$，乘积的指数是 $k+m$。所以配对的依据始终是\textbf{指数相加}；“下标相加”只有在两个级数都从 $z^0$ 开始、指数恰好等于下标时才碰巧一样。

\subhdr{1.2\quad 无限和的情形：定义与一个不能省的条件}

对两个幂级数

\[ \sum_{k=0}^{\infty}a_kz^k,\qquad \sum_{m=0}^{\infty}b_mz^m , \]

定义它们的 \textbf{Cauchy 乘积}为

\[ \sum_{n=0}^{\infty}c_nz^n,\qquad c_n=\sum_{k=0}^{n}a_kb_{n-k} . \]

这个定义只是“把每一项都乘开、按 $z$ 的幂次归堆”，它\textbf{本身并不保证}乘积级数收敛、更不保证收敛到两个和函数的乘积。要保证这件事，需要 §2 的条件。

\subhdr{1.3\quad 一个直观的检查}

取两个几何级数：

\[ \frac{1}{1-z}\cdot\frac{1}{1-z}=\left(\sum_{k\ge0}z^k\right)\left(\sum_{m\ge0}z^m\right) . \]

$z^n$ 的系数是 $k+m=n$ 的配对个数，即 $n+1$：

\[ c_n=\sum_{k=0}^{n}1\cdot1=n+1 . \]

而左边 $=1/(1-z)^2$，它的展开正是 $\sum_{n\ge0}(n+1)z^n$（第三章讲义例 4）。两边一致。这个例子说明：卷积的直观就是“数有多少对配得上”。

\hdr{§2}{什么时候能放心地相乘}

\subhdr{2.1\quad 两条定理}

\textbf{定理（绝对收敛版）}：若 $\sum|a_n|$ 与 $\sum|b_n|$ 都收敛，则 Cauchy 乘积级数 $\sum c_n$ 绝对收敛，且

\[ \sum_{n=0}^{\infty}c_n=\left(\sum_{k=0}^{\infty}a_k\right)\left(\sum_{m=0}^{\infty}b_m\right) . \]

\textbf{定理（Mertens，1880）}：若 $\sum a_n$ 绝对收敛、$\sum b_n$ 收敛（不必绝对），则乘积级数收敛，且和等于两个和的乘积。

\aside{两个都只条件收敛时，乘积级数\textbf{可能发散}——见 2.3 的反例。所以“能乘”的判据落在绝对收敛上，这是本章所有乘法（包括洛朗级数相乘）的依据。}

\subhdr{2.2\quad 为什么绝对收敛就够了（证明思路）}

记部分和 $A_N=\sum_{k=0}^{N}a_k$、$B_N=\sum_{m=0}^{N}b_m$，并记 $C_N=\sum_{n=0}^{N}c_n$。直接算差的“三角形”部分：

\[ A_NB_N=\sum_{k=0}^{N}\sum_{m=0}^{N}a_kb_m , \]

而 $C_N$ 只收集了 $k+m\le N$ 的那些项。两者的差是一圈“边角料”：

\[ A_NB_N-C_N=\sum_{\substack{k+m>N\\ k\le N,\ m\le N}}a_kb_m . \]

若 $\sum|a_k|$、$\sum|b_m|$ 收敛，就来估计这些边角料。注意 $k+m>N$ 而 $k,m\le N$ 时，必有 $k>N/2$ 或 $m>N/2$（否则 $k+m\le N$）。于是

\[ \left|A_NB_N-C_N\right|\le\Bigl(\sum_{k>N/2}|a_k|\Bigr)\Bigl(\sum_{m}|b_m|\Bigr)+\Bigl(\sum_{k}|a_k|\Bigr)\Bigl(\sum_{m>N/2}|b_m|\Bigr) , \]

两个括号里的“尾巴”都随 $N\to\infty$ 趋于 0。所以 $A_NB_N-C_N\to0$，即 $C_N\to AB$。绝对收敛在这里的作用是：\textbf{把“边角料”压到零}。

\subhdr{2.3\quad 反例：两个条件收敛的级数，乘积级数发散}

取

\[ a_n=b_n=\frac{(-1)^n}{\sqrt{n+1}}\qquad(n=0,1,2,\dots) . \]

两个级数都收敛（交错级数，通项递减趋于 0），但不绝对收敛（$\sum1/\sqrt{n+1}$ 发散）。它们的 Cauchy 乘积系数是

\[ c_n=\sum_{k=0}^{n}\frac{(-1)^k}{\sqrt{k+1}}\cdot\frac{(-1)^{n-k}}{\sqrt{n+1-k}}=(-1)^n\sum_{k=0}^{n}\frac{1}{\sqrt{(k+1)(n+1-k)}} . \]

对每一项用算术—几何平均不等式 $\sqrt{xy}\le\dfrac{x+y}{2}$（$x=k+1$，$y=n+1-k$，$x+y=n+2$）：

\[ \frac{1}{\sqrt{xy}}\ge\frac{2}{x+y}=\frac{2}{n+2} . \]

一共 $n+1$ 项，于是

\[ |c_n|\ge\frac{2(n+1)}{n+2}\ \longrightarrow\ 2\ne0 . \]

通项不趋于 0，所以乘积级数发散。数值上，$n=100$ 时这个和约为 $2.85$，与下界 $1.98$ 相容。

\subhdr{2.4\quad 乘积的收敛半径}

若 $\sum a_nz^n$ 的收敛半径为 $R_1$、$\sum b_nz^n$ 为 $R_2$，则在 $|z|<\min(R_1,R_2)$ 内两个级数都绝对收敛，由定理，乘积级数收敛，所以它的收敛半径

\[ R\ge\min(R_1,R_2) . \]

可能更大：例如 $\dfrac{1}{1-z}\cdot(1-z)=1$，乘积的半径是 $\infty$。所以写答案时不要直接把 $\min$ 当成等号。

\hdr{§3}{洛朗级数相乘：指数是负的}

\subhdr{3.1\quad 规则不变，只是把 $m$ 换成 $-m$}

设

\[ f(z)=\sum_{k=-\infty}^{\infty}a_kz^k,\qquad g(z)=\sum_{m=-\infty}^{\infty}b_mz^m \]

在同一个环域内绝对收敛。乘积仍是逐项乘开、按指数归堆：$a_kz^k$ 与 $b_mz^m$ 的乘积指数是 $k+m$，所以

\[ f(z)g(z)=\sum_{n=-\infty}^{\infty}c_nz^n,\qquad c_n=\sum_{k=-\infty}^{\infty}a_kb_{n-k} . \]

这就是双边级数的 Cauchy 乘积，形式和单边的一模一样；区别只在于 $n$ 现在可以取负值。

\textbf{最常见的写法}：把一个因子写成 $z^{-m}$ 的形式，于是“第一项指数 $k$、第二项指数 $-m$”相加：

\[ k+(-m)=n\quad\Longleftrightarrow\quad k-m=n . \]

这正是第三章讲义例 9 里那句“$z^n$ 的系数来自 $k-m=n$ 的那些项”。

\subhdr{3.2\quad 换成 $w=1/z$ 的视角（最不容易错）}

把带负幂的因子看成一个新变量 $w=1/z$ 的幂级数：

\[ e^{-tw/2}=\sum_{m=0}^{\infty}\frac{1}{m!}\left(-\frac t2\right)^mw^m . \]

先对 $z$ 与 $w$ 两个\textbf{独立}变量做标准的 Cauchy 乘积（系数按 $k+m=n$ 配），最后再把 $w$ 换成 $1/z$。这样全程用的都是单边公式，不会把 $k-m$ 与 $k+m$ 搞混。

\subhdr{3.3\quad 收敛域与合法性}

两个洛朗级数在环域 $r_1<|z-z_0|<R_1$ 与 $r_2<|z-z_0|<R_2$ 内绝对收敛，则乘积在交集

\[ \max(r_1,r_2)<|z-z_0|<\min(R_1,R_2) \]

内绝对收敛，并且可以任意重排（绝对收敛级数的重排不变性，第三章讲义 §1.3）。如果交集是空的，就不能直接相乘，要分段处理。

\aside{重排合法性是整件事的根：把 $\sum_{k}\sum_{m}$ 的二维和按 $k+m=n$ 归堆，本质是一次重排。没有绝对收敛，重排可以改变和，甚至破坏收敛（§2.3）。}

\hdr{§4}{三条可操作的算法}

\subhdr{4.1\quad 算法一：直接卷积（列表对角线）}

把两个级数的系数各排一行，$c_n$ 就是“下标之和等于 $n$”的那条对角线上的乘积之和。手算时画一张小表最快：

\begin{center}
\begin{tabular}{c|cccc}
 & $b_0$ & $b_1$ & $b_2$ & $\cdots$ \\ \hline
$a_0$ & $c_0$ & $c_1$ & $c_2$ & \\
$a_1$ & $c_1$ & $c_2$ & & \\
$a_2$ & $c_2$ & & & \\
$\vdots$ & & & &
\end{tabular}
\end{center}

表格的对角线（右上—左下方向）加起来就是 $c_0,c_1,c_2,\dots$。对双边级数，行头列头换成负下标，规则不变（指数相加）。

\subhdr{4.2\quad 算法二：换变量，把双边化单边}

遇到 $z^{-m}$，就令 $w=1/z$，先按单边公式乘，最后把 $w$ 换回 $1/z$。遇到中心不在原点，就令 $u=z-z_0$。换变量的好处是所有配对都回到 $k+m=n$ 的标准形（§3.2）。

\subhdr{4.3\quad 算法三：先化简再乘}

有些乘积根本不必硬卷积：

\begin{itemize}
\item $\sin z\cos z=\tfrac12\sin 2z$（三角恒等式）；
\item $e^{az}e^{bz}=e^{(a+b)z}$（指数律）；
\item $\dfrac{1}{1-z}\cdot\dfrac{1}{1-2z}$ 先部分分式，再分别展开相加。
\end{itemize}

能化简就先化简，卷积只用来处理“两个函数真的不同类”的情形。

\hdr{§5}{五道例题}

\subhdr{例 1\quad $(\sum z^n)^2$}

\[ \left(\sum_{k=0}^{\infty}z^k\right)\left(\sum_{m=0}^{\infty}z^m\right)=\sum_{n=0}^{\infty}\left(\sum_{k=0}^{n}1\right)z^n=\sum_{n=0}^{\infty}(n+1)z^n\qquad(|z|<1) . \]

\textbf{检验}：左边 $=\dfrac{1}{(1-z)^2}$，右边正是它的展开。这道题说明卷积的系数可以“数”出来。

\subhdr{例 2\quad $\dfrac{e^z}{1-z}$（课件 p31 作业第 2 题）}

\[ e^z=\sum_{k\ge0}\frac{z^k}{k!},\qquad \frac{1}{1-z}=\sum_{m\ge0}z^m , \]

两项都在 $|z|<1$ 内绝对收敛。$z^n$ 的系数：

\[ c_n=\sum_{k=0}^{n}\frac{1}{k!}\cdot1=\sum_{k=0}^{n}\frac{1}{k!} . \]

所以

\[ \frac{e^z}{1-z}=\sum_{n=0}^{\infty}\left(\sum_{k=0}^{n}\frac{1}{k!}\right)z^n\qquad(|z|<1) . \]

\textbf{检验}：$z=0$ 给 $1$；一次项系数 $1+1=2$，而 $\left(\dfrac{e^z}{1-z}\right)'_{z=0}=e^0\cdot1+1\cdot e^0=2$，一致。

\subhdr{例 3\quad $\sin z\cos z$：卷积与恒等式对撞}

硬卷积：$\sin z$ 只有奇次项、$\cos z$ 只有偶次项，乘积只含奇次项。$z^{2N+1}$ 的系数是

\[ \sum_{k=0}^{N}\frac{(-1)^k}{(2k+1)!}\cdot\frac{(-1)^{N-k}}{(2N-2k)!}=\frac{(-1)^N\,2^{2N}}{(2N+1)!} . \]

最后这个等式可以逐项验证：$N=1$ 时左边 $=\dfrac{1}{1!\,2!}+\dfrac{1}{3!\,0!}=\dfrac23$，右边 $=\dfrac{4}{3!}=\dfrac23$；$N=2$ 时左边 $=\dfrac{1}{24}+\dfrac{1}{12}+\dfrac{1}{120}=\dfrac{2}{15}$，右边 $=\dfrac{16}{5!}=\dfrac{2}{15}$。于是

\[ \sin z\cos z=\sum_{N=0}^{\infty}\frac{(-1)^N2^{2N}}{(2N+1)!}z^{2N+1}=\frac12\sin 2z . \]

\textbf{结论}：卷积能算出正确的系数，但用恒等式一步就到。两条路给出同一结果，正好互相校验。

\subhdr{例 4\quad $e^{z}\,e^{-1/z}$：双边相乘}

\[ e^{z}=\sum_{k\ge0}\frac{z^k}{k!},\qquad e^{-1/z}=\sum_{m\ge0}\frac{(-1)^m}{m!}z^{-m} . \]

两个级数在 $0<|z|<\infty$ 上都绝对收敛。按指数相加，$z^n$ 的系数（$n$ 可以是负的）来自 $k-m=n$：

\[ c_n=\sum_{m=0}^{\infty}\frac{(-1)^m}{m!\,(m+n)!}\qquad(n\ge0) . \]

这正是 $J_n(2)$ 的级数（因为 $t/2=1$）：

\[ J_n(2)=\sum_{k=0}^{\infty}\frac{(-1)^k}{k!\,(n+k)!} . \]

\textbf{数值检验}：$n=0$ 时 $c_0=\sum_k(-1)^k/(k!)^2=0.2238908$，而 $J_0(2)=0.2238908$；$n=1$ 时 $c_1=0.5767248=J_1(2)$。

\subhdr{例 5\quad 贝塞尔生成函数（课件 p45--p47 的间接法）}

要展开

\[ f(z)=e^{\frac t2\left(z-z^{-1}\right)} . \]

把指数拆成两项，各自展开：

\[ e^{\frac t2z}=\sum_{k=0}^{\infty}\frac{1}{k!}\left(\frac t2\right)^kz^k,\qquad e^{-\frac{t}{2z}}=\sum_{m=0}^{\infty}\frac{1}{m!}\left(-\frac t2\right)^mz^{-m} . \]

两者在 $0<|z|<\infty$ 绝对收敛，可以相乘。按 $k+(-m)=n$ 配对，即 $m=k-n$；对 $n\ge0$，$k$ 从 $n$ 取到 $\infty$：

\[ c_n=\sum_{k=n}^{\infty}\frac{1}{k!}\left(\frac t2\right)^k\cdot\frac{(-1)^{k-n}}{(k-n)!}\left(\frac t2\right)^{k-n} . \]

把 $k$ 换成 $m=k-n$（$m$ 从 $0$ 到 $\infty$）：

\[ c_n=\sum_{m=0}^{\infty}\frac{(-1)^m}{m!\,(m+n)!}\left(\frac t2\right)^{2m+n}=J_n(t) . \]

$n<0$ 时同理得到 $J_n(t)$ 的同一公式（对整数 $n$ 有 $J_{-n}=(-1)^nJ_n$）。所以

\[ e^{\frac t2\left(z-z^{-1}\right)}=\sum_{n=-\infty}^{\infty}J_n(t)\,z^n\qquad(0<|z|<\infty) . \]

\textbf{直接法对照}：取 $C:|z|=1$、$z=e^{i\theta}$，由洛朗系数公式

\[ a_n=\frac{1}{2\pi}\int_0^{2\pi}e^{\frac t2\left(e^{i\theta}-e^{-i\theta}\right)}e^{-in\theta}\dd\theta=\frac{1}{2\pi}\int_0^{2\pi}e^{i\left(t\sin\theta-n\theta\right)}\dd\theta , \]

因为 $\dfrac{e^{i\theta}-e^{-i\theta}}{2}=i\sin\theta$。这就是 $J_n(t)$ 的积分表示，与间接法同一组系数。

\textbf{复核}：$t=0$ 时 $f\equiv1$，右边只剩 $J_0(0)=1$，其余 $J_n(0)=0$。

\hdr{§6}{易错点}

\gapitem{一}{双边级数错用 $k+m=n$。第二项的指数是 $-m$，配对是 $k-m=n$；记不住就先换成 $w=1/z$（§3.2）。}

\gapitem{二}{拿条件收敛的级数硬乘。两个条件收敛的乘积级数可能发散（§2.3），要乘先看绝对收敛，或套 Mertens 条件。}

\gapitem{三}{漏掉负指数项。双边相乘时 $k$ 与 $m$ 都取遍整数，只取 $k,m\ge0$ 会漏掉一半。}

\gapitem{四}{把 $R\ge\min(R_1,R_2)$ 写成等号。$1/(1-z)$ 乘 $(1-z)$ 的半径是 $\infty$，不是 $1$。}

\gapitem{五}{常数项写错。$c_0=a_0b_0$ 永远对，先用它检查一遍再往下算。}

\gapitem{六}{忘记符号。$e^{-t/(2z)}$ 的展开带 $(-1)^m$，漏掉它 $J_n$ 的交替符号就全错了。}

\hdr{§7}{六道自检题（答案在后面）}

\begin{enumerate}
\item $\left(\sum_{k\ge0}\dfrac{z^k}{k!}\right)^2$ 的 $z^n$ 系数是多少？整个级数等于什么函数？
\item $\dfrac{1}{(1-z)^3}$ 的展开是什么？（提示：两个几何级数相乘，或对 $1/(1-z)^2$ 求导。）
\item $e^{z}\cdot e^{-z}$ 的卷积系数 $c_n$ 是多少（$n\ge1$）？为什么能一眼看出来？
\item $\left(\sum_{n\ge0}z^n\right)\left(\sum_{n\ge0}(-1)^nz^n\right)$ 等于什么？卷积里发生了什么？
\item $\dfrac{1}{(1-z)(1-2z)}$ 的展开是什么？（先部分分式，再用卷积对照。）
\item 贝塞尔函数 $J_1(t)$ 的展开前两项是什么？
\end{enumerate}

\textbf{答案}

\begin{enumerate}
\item $c_n=\sum_{k=0}^{n}\dfrac{1}{k!(n-k)!}=\dfrac{2^n}{n!}$，整个级数 $=e^{2z}$。
\item 由 $\sum(n+1)z^n$ 与 $\sum z^m$ 卷积：$c_n=\sum_{k=0}^{n}(k+1)=\dfrac{(n+1)(n+2)}{2}$，即 $\sum\binom{n+2}{2}z^n$。
\item $c_n=\sum_{k=0}^{n}\dfrac{(-1)^k}{k!}\cdot\dfrac{(-1)^{n-k}}{(n-k)!}=\dfrac{(1-1)^n}{n!}=0$（$n\ge1$），$c_0=1$：乘积是常数 1。
\item 卷积 $c_n=\sum_{k=0}^{n}(-1)^{n-k}$ 在 $n$ 偶时为 $1$、$n$ 奇时为 $0$，所以结果是 $\sum z^{2n}=\dfrac{1}{1-z^2}$，与 $1/(1-z)\cdot1/(1+z)$ 一致。
\item 部分分式 $\dfrac{1}{(1-z)(1-2z)}=\dfrac{2}{1-2z}-\dfrac{1}{1-z}$，展开为 $\sum(2^{n+1}-1)z^n$；卷积算 $c_n=\sum_{k=0}^{n}2^k=2^{n+1}-1$，一致。
\item $J_1(t)=\dfrac t2-\dfrac{t^3}{16}+\cdots$（由 $m=0,1$ 两项）。
\end{enumerate}

\hdr{§8}{与课件的对照和一句话总结}

\begin{minipage}{\linewidth}
\small
\tabtop
\begin{tabular}{@{}ll@{}}
\thead{课件页} & \thead{本专题对应处} \\
\end{tabular}
\tabmid
\begin{tabular}{@{}l p{0.72\linewidth}@{}}
p19 & 幂级数的乘法运算；本专题 §1、§2 \\
p20 & 代换运算；本专题 §3.2、§4.2 \\
p45--p47 & 贝塞尔生成函数（直接法、间接法）；本专题 §5 例 5 \\
\end{tabular}
\tabbot
\end{minipage}

\flowerrule

\textbf{一句话}：级数相乘就是\textbf{指数相加}——单边是 $k+m=n$，双边是 $k-m=n$；能这样做的根据是\textbf{绝对收敛}，它保证按指数归堆这一次重排不改变和。

\end{document}